\chapter{Criterio de tabla madre de demos y claims}

\section{Tabla madre}
\begin{definition}
Se llama tabla madre de demos y claims al dispositivo curatorial central mediante el cual el Tomo III organiza, en una sola capa de lectura, las demos seleccionadas, los claims derivados y sus relaciones metodológicas fundamentales.
\end{definition}

\begin{observation}
La tabla madre no sustituye ni a la demo de origen ni a la ficha individual del claim. Su función es ofrecer una vista disciplinada del conjunto, haciendo visible procedencia, jerarquía curatorial, estado y conexiones internas.
\end{observation}

\section{Condición mínima de entrada}
\begin{definition}
Una pieza puede entrar en la tabla madre cuando:
\begin{enumerate}
  \item es una demo seleccionada o un claim ya suficientemente traducido;
  \item conserva trazabilidad suficiente hacia su origen;
  \item posee estado editorial legible;
  \item puede relacionarse con otras piezas sin ambigüedad inmediata de función.
\end{enumerate}
\end{definition}

\begin{observation}
La entrada en la tabla madre no implica igualdad de estatus entre todos los elementos. Solo implica que ya son legibles dentro de una misma cartografía metodológica.
\end{observation}

\section{Campos mínimos}
\begin{definition}
La tabla madre debe permitir, al menos, los siguientes campos:
\begin{enumerate}
  \item identificador;
  \item tipo de pieza (demo o claim);
  \item fuente u origen;
  \item dependencia canónica;
  \item estado editorial;
  \item prioridad curatorial;
  \item observación crítica, si existe;
  \item vínculo con otras piezas relacionadas.
\end{enumerate}
\end{definition}

\begin{observation}
Estos campos fijan una primera gramática común de inventario curado. No agotan toda estructura futura de la tabla madre.
\end{observation}

\section{Tabla madre y continuidad}
\begin{definition}
La tabla madre debe hacer visible, de forma suficiente, la continuidad entre:
\[
\text{admisión metodológica}
\;\longrightarrow\;
\text{demo seleccionada}
\;\longrightarrow\;
\text{claim traducido}
\;\longrightarrow\;
\text{estado curatorial}.
\]
\end{definition}

\begin{observation}
Esta continuidad no obliga a que toda demo termine en claim ni a que todo claim alcance prioridad alta. Solo organiza el recorrido posible de cada pieza dentro del tomo.
\end{observation}

\section{Condición mínima de no entrada}
\begin{definition}
Una pieza no debe entrar en la tabla madre cuando su origen sigue siendo opaco, su estatus es todavía inestable o su relación con el resto del corpus no puede formularse con suficiente sobriedad.
\end{definition}

\begin{observation}
Esta condición protege la tabla madre frente a inflación de inventario y evita que funcione como depósito indiferenciado.
\end{observation}

\section{Cláusula de alcance}
\begin{observation}
El criterio de tabla madre fija una primera cartografía común de demos y claims dentro del Tomo III. No constituye todavía una versión pública definitiva ni una clausura exhaustiva del mapa metodológico del corpus.
\end{observation}

